\documentclass[11pt,twoside]{article}
\usepackage[paperwidth=17.00cm,paperheight=25.00cm,top=2.20cm,inner=2.50cm,outer=1.65cm,bottom=2.20cm]{geometry}
\usepackage{fontspec}
\usepackage{polyglossia}
\usepackage[style=ascph,backend=biber,url=true]{biblatex}
\addbibresource{bibliography.bib}
\usepackage[breaklinks=true,hidelinks]{hyperref}
\setmainlanguage{spanish}
\setmainfont{Alt Standard}[RawFeature=lnum]
\usepackage{microtype}
\sloppy
%
\usepackage{titleps}
\newcommand{\head}[1]{\def\mihead{#1}\thispagestyle{empty}} \def\mihead{} \newpagestyle{editio}{\sethead[\thepage][\ifx\mihead\empty\else\mihead\fi][]{}{\ifx\mihead\empty\else\mihead\fi}{\thepage}} \pagestyle{editio}
%
\title{\huge \textbf{ALIA STRVCTVRA CITATIONIS PHILOLOGICÆ}\\ \textsc{\huge\textbf{--ASCPh--}}}
\date{}
\author{\huge\textsc{\textbf{Ioannes Bracmanvs}}}
%
\begin{document}
		{\pagenumbering{gobble}
		\maketitle
	}
	\newpage
%
\section*{\centering\textsc{Praefatio}}
%
\textsc{Aqvellos} posibles lectores de esta propuesta podrán intuir, por el título, que esta estructura fue diseñada únicamente para el contexto filológico, y los invadirá la duda de ¿qué sentido tiene «otra» estructura si ya existen las canónicas, excelentes para el uso académico, como APA o MLA? Ciertamente, es una duda genuina y atendible.
%
\section*{\centering\textsc{Codicvm traditio}}
%
%
\section*{\centering\textsc{Editoris primatvs}}
%
%
\section*{\centering\textsc{Forma bibliographica}}
%
%
\section*{\centering\textsc{Citationvm formæ}}
%
%
\section*{\centering\textsc{Strvctvra et exempla}}
Para libros con textos modernos = \textbf{Autor, A. (Fecha). \emph{Título} (Ed./Tr.= Nombre, N. y Nombre, N.). (Vol. X, T. Y/Fasc. Y, Z\textsuperscript{ed.}). Editorial. <Enlace>. (= \emph{Sigla}).}\par
%
Para libros con textos antiguos = \textbf{Editor, E. (Fecha). \emph{Título} (\emph{Au.}= Nombre, N. \& Nombre, N.). (Vol. X, T. Y/Fasc. Y, Z\textsuperscript{ed.}). Editorial. <Enlace>. (= \emph{Sigla}).}\par
%
Con texto de revistas = \textbf{Autor, A. (Fecha). ``Título''. Revista, Nº, páginas. Editorial. <Enlace>.}\par
%
Para capítulos o artículos en libros o compilaciones = \textbf{Autor, A. (Fecha). ``Título''. \emph{Título del libro o compilación} (Serie o «colección», Nº), Nº págs. Editorial. <Enlace>.}\par
%
Con texto o material desde internet = \textbf{Autor, A. (Fecha/S. F.). ``Nombre''. <Enlace>. (= \emph{Sigla}).}\par
%
\vskip11pt
\textbf{Cita narrativa =} «Así como señala \textcite{west2017}, es preciso...»/ \textbf{Cita literal =} «``... Odysseam non in Ionia compositam esse sed in Graecia'' \parencite[VII]{west2017}»/ \textbf{Cita literal en la prosa =} «... ciertamente porque, en palabras de West, ``Odysseam non in Ionia compositam esse sed in Graecia'' \prosecite[VII]{west2017}».\par
%
\noindent- - - - - - - - - - - - -\par
%
\nocite{*}
\printbibliography[heading=none]
\end{document}